\documentclass[12pt]{article}
% Préambule commun à tous les documents extraits de « La Revue CAPES »
\usepackage[T1]{fontenc}
\usepackage[utf8]{inputenc}
\usepackage{lmodern}
\usepackage{amsmath,amssymb,amsthm,amsfonts,mathrsfs}
\usepackage{setspace}
\usepackage{listings}
\usepackage{pgfplots}
\pgfplotsset{compat=1.18}
\usepackage{longtable}
\usepackage{adjustbox}
\usepackage{lettrine}
\usepackage{tkz-tab}
\usepackage{changepage}
\usepackage{pdflscape}
\usepackage{graphicx}
\usepackage{tikz}
\usepackage{xcolor}
\usepackage[a4paper,top=2cm,bottom=2.5cm,left=2.5cm,right=2cm]{geometry}
\usepackage{fancyhdr}
\usepackage{draftwatermark}
\SetWatermarkText{La RevueCapes}
\SetWatermarkLightness{0.9}
\SetWatermarkScale{2}
\usepackage{hyperref}
\hypersetup{colorlinks=true,linkcolor=blue!50!black,urlcolor=blue!50!black}

\definecolor{revue}{RGB}{20,60,120}

\pagestyle{fancy}
\fancyhf{}
\renewcommand{\headrulewidth}{0.4pt}
\setlength{\headheight}{15pt}

% \entete{Matière}{Titre}{Type de document}
\newcommand{\entete}[3]{%
  \fancyhead[L]{\small\textsc{La Revue CAPES} -- #1}%
  \fancyhead[R]{\small #2}%
  \fancyfoot[C]{\thepage}%
  \thispagestyle{plain}%
  \begin{center}
    {\color{revue}\rule{\textwidth}{1.2pt}}\\[4pt]
    {\small\textsc{Concours directs CAP-CEG / CAPES -- Burkina Faso}}\\[6pt]
    {\LARGE\bfseries\color{revue} #2}\\[6pt]
    {\large\itshape #1 -- #3}\\[2pt]
    {\color{revue}\rule{\textwidth}{1.2pt}}
  \end{center}
  \vspace{0.5cm}%
}

% Encadré pour les remarques ajoutées dans les corrigés
\newcommand{\remarque}[1]{\par\medskip\noindent\fcolorbox{revue}{revue!6}{\parbox{\dimexpr\linewidth-2\fboxsep-2\fboxrule}{\textbf{\color{revue}Remarque.} #1}}\par\medskip}


\begin{document}
\entete{Mathématiques}{Ancien sujet CAPES -- Session 2019}{Énoncé}
\begin{spacing}{1.5}

\textbf{Exercice 1 : (8pts)}

On considère la matrice carrée d'ordre 3 définie par  : 

\[M=\begin{pmatrix}
11 & -5 & 5 \\ 
-5 & 3 & -3 \\ 
5 & -3 & 3
\end{pmatrix}\]

1. Déterminer le polynôme caractéristique de $M$, ainsi que son spectre.

2. Quels sont les sous espaces propres associés aux différentes valeurs propres?

3. $M$ est-elle diagonalisable ? Si oui, la diagonaliser. On précisera la base $B$ et les matrices $D$ et $P$ de cette base.

4. En déduire $M^{n}$ $\forall n\in \mathbb{N}$.

\textbf{Exercice 2 : (6pts)}

On considère la fonction $2\pi$ périodique impaire définie par : 

$h(x)=x(\pi-x); \forall {x}\in [0,\pi]$.

1. Représenter graphiquement cette fonction sur $[-2\pi,2\pi]$.

2. Calculer les coefficients de Fourier de $h$.

3. En déduire le développement en série de Fourier de $h$.

4. Calculer alors les sommes(exactes) des séries numériques 

$S_1 = \sum^{+\infty}_{n=0}\frac{(-1)^n}{(2n+1)^3}$ et $S_2 = \sum^{+\infty}_{n=0}\frac{1}{(2n+1)^6}$

\textbf{Exercice 3 : (5pts)}

Un industriel doit vérifier l'état de marche de ses machines et en remplacer certaines le cas échéant. D'après des statistiques précédentes, il évalue à 30% la probabilité pour une machine de tomber en panne en cinq ans. Parmi ces dernières, la probabilité de devenir hors d'usage suite à une panne plus grave est évaluée à 75%. Cette probabilité est de 40% pour une machine n'ayant jamais eu de panne.

1. Quelle est la probabilité pour une machine donnée de plus de cinq ans d'être hors d'usage?

2. Quelle est la probabilité pour une machine hors d'usage de n'avoir jamais eu de panne au paravent?

3. Soit $X$ la variable aléatoire $\ll$ {la nombre de machine qui tombe en panne au bout de cinq ans, parmi 10 machines choisies au hasard}$\gg$. Quelle est la probabilité de $X$(on donnera le type de la loi et la formule de calcul), son espérance, sa variance et son écart-type?

\end{spacing}
\end{document}
